\chapter{Estudio del problema}
\label{Estudio del problema}

Este capítulo sitúa el trabajo. Comienza por el contexto industrial que motiva la tarea y por la
plataforma robótica elegida, y justifica por qué todo el desarrollo se hace en simulación. Revisa
después el estado de la cuestión en las áreas que el trabajo cruza —aprendizaje por imitación,
políticas entrenadas desde cero, modelos visión-lenguaje-acción, control de humanoides, simulación,
teleoperación y reconstrucción de escenas—, delimita a continuación el problema concreto que se
aborda, con sus limitaciones, y termina recogiendo la normativa que afectaría a un despliegue real
del sistema.


\section{El contexto del problema}

\subsection{La operación de válvulas en instalaciones de proceso}

Una instalación de proceso —una refinería, una planta de tratamiento de gas, una plataforma de
extracción— es, vista de cerca, una red de tuberías interrumpida cada pocos metros por una válvula.
Las de compuerta y las de globo accionadas por volante son las más numerosas: son baratas, robustas
y no necesitan más energía que la del brazo de quien las gira. Una válvula de compuerta abre y cierra
desplazando un disco perpendicular al flujo mediante un husillo roscado, de modo que pasar de cerrada
a completamente abierta exige varias vueltas completas del volante; el modelo CAD empleado en este
trabajo, tomado de una válvula real, autoriza seis vueltas y media de recorrido. El par necesario
crece con el diámetro nominal y con la presión de la línea, y en una válvula que lleva meses sin
moverse hay que sumar el de arranque.

Pese a la extensión de la instrumentación, una fracción muy grande de esas válvulas sigue siendo
manual y va a seguir siéndolo: motorizarlas todas es prohibitivo, y muchas solo se accionan en
paradas, arranques o emergencias. Su operación recae en personal de campo que recorre la planta
siguiendo hojas de ruta, y tiene tres características que la hacen candidata natural a la
automatización. Es \textbf{repetitiva}: el gesto es el mismo válvula tras válvula. Es
\textbf{físicamente exigente}: volantes a ras de suelo o por encima de la cabeza, pares de arranque
altos, posturas forzadas. Y se realiza en \textbf{entornos hostiles}: temperaturas extremas, ruido y,
sobre todo, zonas clasificadas como atmósferas potencialmente explosivas, en las que la mera
presencia de una persona es un riesgo que hay que gestionar y en las que, durante un incidente, la
operación manual es lo último que se quiere tener que hacer.

No es casual que girar una válvula fuera una de las pruebas del \textit{DARPA Robotics Challenge}
\cite{krotkov2017drc}, el concurso que en 2015 puso a los humanoides a realizar tareas de
intervención en un escenario de accidente industrial, ni que las iniciativas de robótica para el
sector —el reto ARGOS promovido por Total, o los despliegues de cuadrúpedos de inspección en
plataformas marinas— la cuenten entre sus capacidades objetivo. Es una tarea con un criterio de
éxito objetivo —la válvula abre o no abre—, bien delimitada y con valor directo. Y es, al mismo
tiempo, una tarea que resiste mal el enfoque clásico de programar la trayectoria: la válvula puede
estar montada en cualquier orientación, el volante exige un agarre firme y un giro sostenido, y el
par de reacción se transmite por los brazos hasta la base del robot.

\subsection{Robots humanoides como plataforma de intervención}

La pregunta obvia es por qué un humanoide y no un manipulador sobre base móvil, que es una
tecnología madura. La respuesta está en la instalación, no en el robot. Una planta de proceso está
diseñada para el cuerpo humano: los pasillos tienen la anchura de una persona, los accesos son
escaleras y pasos elevados de peldaños, los volantes están a la altura y a la distancia a la que un
operario los alcanza de pie, y su diámetro es el que una mano agarra. Un robot con la morfología del
operario hereda ese diseño sin reformar nada; cualquier otra morfología obliga a adaptar la planta o
a renunciar a parte de ella.

Lo que ha cambiado en los últimos tres años no es ese argumento, que es antiguo, sino el precio de la
plataforma. Los humanoides de investigación de la década pasada eran prototipos únicos con
presupuestos de siete cifras; hoy varios fabricantes venden humanoides completos por el precio de un
vehículo industrial. El Unitree G1 \cite{unitree_g1}, empleado en este trabajo, es representativo de
esa generación: un humanoide de tamaño reducido —del orden de 1,3\,m y 35\,kg—, con veintinueve
grados de libertad en el cuerpo y, en la variante utilizada, manos de tres dedos con siete
articulaciones actuadas cada una, hasta un total de cuarenta y tres. Es además la plataforma para la
que NVIDIA distribuye su pila de aprendizaje para humanoides, lo que a efectos de este trabajo
significa que existen un modelo de simulación mantenido, una política de equilibrio de cuerpo
completo ya entrenada y una descripción del \textit{embodiment} para el modelo fundacional. Elegir
otro robot habría supuesto reconstruir esas tres cosas.

Su tamaño es también una restricción. Un G1 de pie alcanza con comodidad un volante frontal a unos
43\,cm del pelvis y uno cenital a unos 30\,cm, distancias que se han medido con el operador dentro
del sistema de teleoperación y que fijan dónde puede estar la válvula. Un humanoide de talla adulta
tendría más alcance y más par; también más masa que controlar cuando el volante devuelve la
reacción.

\subsection{Por qué el trabajo se realiza en simulación}
\label{sec:por-que-simulacion}

El trabajo se realiza íntegramente en simulación por una razón metodológica y por dos prácticas.

La metodológica es la que dimensiona el experimento. Comparar dos políticas con rigor exige
evaluarlas muchas veces —en este trabajo, doscientas tiradas por política en la comparación
principal y seiscientas más en la curva de eficiencia— y, sobre todo, exige evaluarlas sobre las
\textbf{mismas} condiciones iniciales, de modo que cada episodio sea un par y el contraste
estadístico pueda aprovecharlo. Sobre un robot físico, ni lo uno ni lo otro es viable en el plazo de
un trabajo de fin de máster: reproducir la pose de una válvula al milímetro cientos de veces no es
posible, y cada tirada cuesta minutos de preparación y expone al robot a una caída. En simulación,
una semilla fija garantiza que la tirada $i$ de una política ve exactamente la escena que vio la
tirada $i$ de la otra, y mil episodios son una noche de máquina.

A pesar de haber dispuesto de un G1 físico durante el desarrollo, recoger demostraciones por teleoperación sobre el robot real habría requerido una infraestructura de seguridad —arnés, zona acotada, un segundo operador— que no estaba
al alcance del proyecto.

Conviene ser explícito sobre lo que se pierde. Primero, la dinámica de contacto: el agarre del
volante y el par que devuelve son, en el simulador, lo que el motor de física y los parámetros de la
articulación dicen que son, no lo que hace un volante de fundición. Segundo, el ruido de los sensores
y las latencias de la plataforma real. Tercero, y más importante, que la política de equilibrio que
sostiene al robot está aquí liberada de parte de su trabajo: la base del robot se fija durante la
manipulación, una decisión que se justifica en el capítulo \ref{Solución} y que en un despliegue
real no existiría.

Se ha intentado mitigar esa distancia por tres vías. La válvula procede de un modelo CAD de una
válvula real, no de una primitiva geométrica. Las imágenes que ven las políticas no muestran un
fondo sintético, sino la reconstrucción fotorrealista de una oficina real obtenida a partir de vídeo
(sección \ref{sec:gs-estado} y anexo \ref{Reconstrucción fotorrealista}). Y la evaluación no se
limita a la tasa de éxito: mide también la latencia de inferencia en lazo cerrado, que es la
magnitud que decidiría si la política cabe en el robot.


\section{El estado de la cuestión}

\subsection{Aprendizaje por imitación}

El aprendizaje por imitación es la familia de técnicas que obtienen una política de control a
partir de demostraciones, en lugar de a partir de una función de recompensa. Su forma más simple, la
clonación de comportamiento (\textit{behavior cloning}), reduce el problema a uno de aprendizaje
supervisado: se recogen pares observación--acción de un experto y se ajusta un modelo que prediga la
segunda a partir de la primera. Es una idea antigua. ALVINN \cite{pomerleau1988alvinn} condujo un
vehículo en 1988 con una red de tres capas que asociaba la imagen de una cámara a un ángulo de
dirección, entrenada con las maniobras de un conductor humano.

Su limitación fundamental también se conoce desde entonces, y tiene nombre: la deriva de
distribución. La política se entrena sobre los estados que el experto visitó, pero al ejecutarse
comete pequeños errores que la llevan a estados ligeramente distintos, sobre los que nunca vio una
demostración y en los que sus errores son mayores; el proceso se realimenta. Ross \textit{et al.}
\cite{ross2011dagger} formalizaron el problema —en el peor caso, el coste acumulado de la clonación
ingenua crece con el cuadrado de la longitud del episodio— y propusieron DAgger como respuesta:
alternar la ejecución de la política con el etiquetado por el experto de los estados que esta
visita, de modo que el conjunto de entrenamiento acabe cubriendo la distribución que la propia
política induce. DAgger exige un experto disponible en línea, lo que en la teleoperación de un
humanoide es costoso, y por eso las aproximaciones actuales atacan la deriva por otro lado:
prediciendo bloques de acciones en lugar de acciones aisladas, lo que reduce el número de decisiones
en las que la política puede equivocarse y da a cada una un contexto temporal más largo; y, en los
modelos fundacionales, apoyándose en un preentrenamiento que ha visto muchos más estados que las
demostraciones de la tarea.

Se fija aquí el vocabulario que emplea el resto de la memoria. Una \textbf{demostración} es un
episodio completo ejecutado por el operador y grabado. Un \textbf{episodio} de evaluación es un
intento de la política, que termina en éxito o por agotamiento del tiempo; cada uno de esos intentos
es una \textbf{tirada} (\textit{rollout}). Una \textbf{política} es la función que asocia
observaciones a acciones. Todos ellos, y los demás términos técnicos, están recogidos en el
glosario.

\subsection{Políticas de manipulación entrenadas desde cero}

Dentro de las arquitecturas específicas —entrenadas exclusivamente sobre las demostraciones de la
tarea, sin preentrenamiento— dos han dominado los últimos tres años.

ACT (\textit{Action Chunking with Transformers}) \cite{zhao2023aloha} nació junto con ALOHA, un
sistema bimanual de bajo coste, con un objetivo declarado: aprender manipulación fina con pocas
demostraciones, del orden de cincuenta, y hardware imperfecto. Su aportación es doble. La primera es
el \textit{chunking}: la política no predice la siguiente acción sino un bloque de acciones futuras
—en el trabajo original, un centenar de pasos a 50\,Hz—, lo que reduce el horizonte efectivo de la
tarea y con él la acumulación de errores. La segunda es modelar las demostraciones con un
autocodificador variacional condicional (CVAE): durante el entrenamiento, un codificador comprime la
secuencia de acciones en una variable latente, y un decodificador basado en \textit{transformer}
genera el bloque a partir de las imágenes, del estado articular y de esa variable; en inferencia, el
latente se fija a cero y la política es determinista. El resultado es una arquitectura compacta, de
decenas de millones de parámetros, rápida y muy extendida, lo que la convierte en el representante
natural del brazo clásico de esta comparación.

Diffusion Policy \cite{chi2023diffusion} es la alternativa generativa. En lugar de regresar el bloque
de acciones, lo obtiene por un proceso de difusión que parte de ruido y lo refina en varios pasos
condicionados por la observación. Su ventaja está en las demostraciones multimodales: cuando hay
varias formas válidas de hacer lo mismo, una regresión aprende la media, que no es ninguna de ellas,
y un modelo generativo puede elegir una. A cambio, cada inferencia requiere varias evaluaciones de
la red.

Se ha elegido ACT y no Diffusion Policy para el brazo clásico por tres razones. Es la arquitectura
más extendida y mejor documentada, con una implementación de referencia mantenida en LeRobot. Su
inferencia es de una sola pasada, lo que hace más limpia la medida de latencia. Y la tarea, tal como
se ha grabado —un único operador, un gesto de agarre fijado de antemano—, no es especialmente
multimodal, de modo que la ventaja principal de la difusión no tendría ocasión de manifestarse.
Nada impide añadir Diffusion Policy como tercer brazo: el \textit{pipeline} de datos y el banco de evaluación
le servirían sin cambios.

\subsection{Modelos visión-lenguaje-acción}

La otra familia parte de una observación: los grandes modelos de visión y lenguaje ya saben mucho
del mundo —qué es una válvula, dónde tiene el volante, qué significa «abrir»— y ese conocimiento
puede transferirse al control si se enseña al modelo a emitir acciones.

La trayectoria es reciente y rápida. RT-1 \cite{brohan2022rt1} demostró en 2022 que un
\textit{transformer} entrenado sobre unos 130\,000 episodios reales, recogidos durante diecisiete
meses con una flota de trece robots, generalizaba a cientos de tareas descritas en lenguaje natural.
RT-2 \cite{brohan2023rt2} dio el paso decisivo: en lugar de entrenar un modelo de control desde cero,
tomó modelos visión-lenguaje preentrenados sobre la web y los ajustó para que produjeran las
acciones como si fueran texto —una secuencia de \textit{tokens} que codifican el desplazamiento del
efector—. Acuñó el término \textit{vision-language-action} y mostró generalización a objetos y
conceptos que no aparecían en los datos robóticos. OpenVLA \cite{kim2024openvla} llevó esa receta a
un modelo abierto de siete mil millones de parámetros, entrenado sobre casi un millón de episodios
de Open X-Embodiment \cite{oxe2023}, y mostró que el ajuste fino con adaptadores de bajo rango basta
para adaptarlo a un robot nuevo. $\pi_0$ \cite{black2024pi0} cambió la forma de generar la acción: en
lugar de \textit{tokens} discretos, un «experto de acción» separado produce bloques continuos de
acciones mediante \textit{flow matching}, lo que permite emitirlas a 50\,Hz con la destreza que
exige, por ejemplo, doblar ropa.

Lo que esta familia promete es generalización —a instancias del objeto no vistas, a variaciones de
la escena, a instrucciones nuevas— con menos demostraciones de la tarea concreta que una
arquitectura entrenada desde cero. Lo que cuesta es cómputo. Los VLA tienen entre uno y dos órdenes
de magnitud más parámetros que ACT, su ajuste fino requiere aceleradoras de gama profesional y su
inferencia tarda decenas de milisegundos donde ACT tarda unidades. Si esa diferencia importa depende
del presupuesto de tiempo del lazo de control, y por eso se mide en el capítulo \ref{Evaluación}.

\subsection{GR00T y el control de humanoides}

GR00T N1 \cite{nvidia2025gr00t} es el modelo fundacional de NVIDIA para humanoides, y su variante
N1.7 es la empleada en este trabajo. Su arquitectura es de dos sistemas, por analogía con la
cognición humana. El «sistema 2» es un modelo visión-lenguaje que procesa las imágenes y la
instrucción y produce una representación de la escena a baja frecuencia. El «sistema 1» es un
\textit{transformer} de difusión que, condicionado por esa representación y por el estado del robot,
genera bloques de acciones continuas mediante \textit{flow matching}, en la línea de $\pi_0$. Ambos
se entrenan de extremo a extremo.

Dos rasgos lo hacen especialmente adecuado a este trabajo. El primero es que está diseñado para
servir a varios \textit{embodiments}: el estado y la acción de cada robot pasan por codificadores y
decodificadores propios —pequeñas redes que se entrenan durante el ajuste fino— mientras el núcleo
del modelo es común. Incorporar el G1 con la configuración de este trabajo consiste en describir su
espacio de estado y de acción, no en modificar el modelo. El segundo es la «pirámide de datos» con
la que se preentrenó: en la base, vídeo humano y de la web; en el medio, datos sintéticos generados
en simulación; en la cúspide, teleoperación real de humanoides. Es un preentrenamiento que ha visto
humanoides manipulando, lo que no puede decirse de los VLA anteriores, entrenados sobre brazos de
sobremesa.

El modelo produce acciones para el tren superior; mantener al humanoide de pie mientras las ejecuta
es un problema distinto, y en este trabajo lo resuelve una capa separada. Es la arquitectura habitual
en manipulación con humanoides: una política de control de cuerpo completo (\textit{whole-body
control}, WBC), entrenada por refuerzo en simulación, gobierna las piernas y mantiene el equilibrio,
y recibe de arriba órdenes de alto nivel —velocidad de avance, altura del pelvis, orientación del
torso—, mientras las poses de las muñecas que marca la política de manipulación se convierten en
ángulos articulares mediante cinemática inversa. La política de imitación aprende, por tanto, sobre
un robot que ya sabe estar de pie. Esa división es lo que hace tratable el problema: el espacio de
acción que aprende la política tiene veintitrés dimensiones, no cuarenta y tres.

\subsection{Entornos de simulación y formatos de conjuntos de datos}

La pila de simulación empleada es la de NVIDIA. Isaac Sim \cite{isaacsim} es el simulador: motor de
física PhysX, render por trazado de rayos y escenas descritas en USD. Isaac Lab
\cite{mittal2023orbit}, nacido como Orbit, es el marco de aprendizaje construido sobre él: aporta la
ejecución paralela en GPU de miles de entornos, la definición modular de una tarea mediante gestores
de observaciones, acciones, recompensas y terminaciones, y los adaptadores para las bibliotecas de
aprendizaje. Isaac Lab Arena \cite{isaaclab_arena} es la capa más reciente: compone escenas de
manipulación a partir de un registro de activos —robots, objetos, fondos— y añade la infraestructura
de teleoperación, grabación de demostraciones, conversión de datos y evaluación de políticas,
incluidas las que se ejecutan en otro proceso.

Se consideraron y descartaron dos alternativas. MuJoCo, por su precisión de contacto y su bajo coste
de cómputo, habría sido una opción razonable para la física, pero no dispone del sistema de
teleoperación en realidad mixta ni del \textit{embodiment} del G1 con su política de equilibrio ya
integrada, y reconstruir ambas cosas no cabía en el plazo. Dentro de la propia pila de NVIDIA se
consideró generar las demostraciones automáticamente con Isaac Lab Mimic, que multiplica un puñado
de demostraciones humanas recombinando sus segmentos, y se rechazó como riesgo de calendario: exige
anotar la estructura de la tarea y no estaba probado sobre este \textit{embodiment}.

En cuanto a los datos, el denominador común es el formato LeRobot \cite{cadene2024lerobot}, de
Hugging Face: cada episodio como filas de un fichero \textit{parquet} con el estado y la acción, las
imágenes como vídeo comprimido, y metadatos que describen las modalidades. Es el formato que
consumen las implementaciones de referencia de ACT y de Diffusion Policy, y el que GR00T adopta con
una extensión —un fichero de modalidades que asigna rangos del vector de estado a partes del
cuerpo— que ACT simplemente ignora. Esa compatibilidad es la que permite que una sola conversión
alimente a los dos métodos. Detrás del formato hay un movimiento más amplio de agregación de datos
robóticos, del que Open X-Embodiment \cite{oxe2023} —más de un millón de trayectorias de veintidós
\textit{embodiments}, aportadas por veintiuna instituciones— es el exponente principal, y que es lo
que ha hecho posibles los VLA abiertos.

\subsection{Teleoperación en realidad mixta}

Las demostraciones hay que grabarlas, y la forma de hacerlo condiciona su calidad. Para brazos de
sobremesa la solución dominante son los dispositivos maestro--esclavo —ALOHA es el ejemplo—, en los
que el operador mueve una réplica pasiva del robot. Para un humanoide no hay réplica que mover: la
captura debe hacerse sobre el cuerpo del operador.

La vía que se ha impuesto es la realidad mixta con dispositivos de consumo. Unas gafas con
seguimiento de cabeza y de manos —o mandos— proporcionan la pose de la cabeza y de las muñecas del
operador varias decenas de veces por segundo; un módulo de \textit{retargeting} traslada esas poses
al cuerpo del robot, que tiene otras proporciones y otras articulaciones; y el vídeo estereoscópico
de las cámaras del robot se devuelve a las gafas, de modo que el operador ve lo que el robot ve y
cierra el lazo con su propia percepción. Open-TeleVision \cite{cheng2024television} es una
referencia reciente de este esquema.

En este trabajo las gafas son unas PICO 4 Ultra \cite{pico4ultra}, un visor autónomo: lleva su
propio procesador y no está unido a ningún ordenador por cable. Eso plantea un problema, porque el
simulador necesita una aceleradora de gama profesional y no puede ejecutarse en el visor. Lo
resuelve NVIDIA CloudXR \cite{nvidia_cloudxr}, la capa de transmisión que separa dónde se renderiza
de dónde se muestra: el simulador genera las dos vistas estereoscópicas en la estación de trabajo,
CloudXR las codifica como vídeo y las envía por la red al visor, y en sentido contrario recibe la
pose de la cabeza y de los mandos varias decenas de veces por segundo y se la entrega al simulador
a través de la interfaz OpenXR. Para el sistema, el visor es un cliente de vídeo con sensores; para
el operador, está dentro de la escena. Los mandos son la fuente de la pose de las muñecas y el
gatillo la orden de cierre de la mano; el capítulo \ref{Solución} documenta el montaje y los
problemas que dio.

Dos propiedades de esta vía importan para lo que viene después. La primera es que el operador forma
parte del sistema: su fatiga, sus hábitos y sus errores quedan grabados, y la política los copia. Por
eso las sesiones de grabación se limitan a veinticinco demostraciones, y por eso el gesto de agarre
se fijó antes de grabar el conjunto definitivo. La segunda es que la latencia entre el movimiento
del operador y la imagen que recibe determina cuánto puede afinar; un fondo costoso de renderizar
la empeora, y ese es el motivo por el que el fondo fotorrealista no se aplica durante la grabación.

\subsection{Realismo visual mediante reconstrucción de escenas}
\label{sec:gs-estado}

Una política que aprende de imágenes aprende también del fondo. Si las imágenes de entrenamiento
muestran un plano gris bajo una luz uniforme, la política no tiene motivo para ignorarlo y, puesta
ante una escena real, verá algo que nunca vio. Los fondos sintéticos tradicionales, modelados a mano,
mitigan el problema a un coste de modelado alto y con un realismo limitado.

La reconstrucción neuronal de escenas ofrece otra vía: capturar un entorno real con una cámara y
obtener una representación que pueda renderizarse desde cualquier punto de vista. Los campos de
radiancia neuronales abrieron la línea, pero su render era demasiado lento para un simulador
interactivo. El \textit{Gaussian Splatting} tridimensional \cite{kerbl2023gaussian} resolvió esa
limitación representando la escena como millones de gaussianas anisótropas explícitas —cada una con
posición, covarianza, opacidad y color dependiente de la dirección de vista— que se optimizan a
partir de las imágenes y se rasterizan en tiempo real con una calidad comparable a la de los mejores
campos de radiancia. NVIDIA lo ha integrado en Omniverse bajo el nombre de NuRec, lo que permite
cargar una de esas reconstrucciones como un activo más de una escena USD y renderizarla desde la
cámara del robot.

En este trabajo el fondo es la reconstrucción de una oficina real. Su papel es acercar las imágenes
de entrenamiento a las de un despliegue, y su coste —unos cuarenta milisegundos por paso de
simulación— es lo que obliga a aplicarlo a posteriori, reproduciendo las demostraciones con el
fondo cargado. Cómo se hace, y el problema que planteó, se desarrollan en el anexo
\ref{Reconstrucción fotorrealista}.


\section{La definición del problema}

El problema que aborda este trabajo cabe en una frase: \textbf{dado un conjunto de demostraciones
teleoperadas de la apertura de una válvula de volante con un humanoide, determinar, en condiciones
idénticas y con el contraste estadístico apropiado, si un modelo fundacional visión-lenguaje-acción
ajustado a la tarea supera a una política de imitación clásica entrenada desde cero, y a qué coste
en datos y en latencia.} Resolverlo exige, antes, construir el sistema que lo hace medible: la
tarea, el \textit{pipeline} de datos y el banco de evaluación.

El argumento que lo motiva es el vacío que deja la literatura. Cada familia acumula resultados por
separado: los VLA publican tasas de éxito sobre sus propios conjuntos de tareas, y las arquitecturas
específicas sobre los suyos. Cuando se comparan entre sí, suele ser dentro del artículo que presenta
uno de los métodos, sobre tareas elegidas por sus autores, con conjuntos de datos que no siempre
coinciden y con una comparación de proporciones independientes como único contraste, cuando hay
alguno. Faltan comparaciones en las que la tarea, los datos, el criterio de éxito y las condiciones
iniciales sean los mismos para ambos métodos, y en las que el contraste aproveche ese
emparejamiento. Sin ellas, la elección del método para un despliegue concreto se toma sobre el
tamaño del modelo y la reputación de la familia, no sobre evidencia.

Para que la comparación sea limpia, el trabajo asume las restricciones siguientes, que son a la vez
sus limitaciones:

\begin{itemize}

	\item \textbf{Una sola tarea.} La apertura de la válvula, en dos disposiciones y con la posición
	aleatorizada. Los resultados describen esa tarea; lo que se transfiere a otras es el banco de
	comparación, no las cifras.

	\item \textbf{Sin validación sobre el robot físico.} Todo ocurre en simulación, con las pérdidas descritas en la
	sección \ref{sec:por-que-simulacion}. La distancia entre el simulador y el robot no se mide.

	\item \textbf{Una única semilla de entrenamiento por configuración.} Cada punto de la curva de
	eficiencia es un solo entrenamiento, de modo que la variación entre puntos mezcla el efecto de
	los datos con la varianza del propio entrenamiento.

	\item \textbf{Un único operador.} Las cien demostraciones las grabó la misma persona. Eso da un
	gesto consistente, que facilita el aprendizaje, y a la vez sesga el conjunto hacia un estilo.

	\item \textbf{Una base fija.} La pelvis del robot se inmoviliza durante la manipulación; la
	política de equilibrio sigue activa, pero no tiene que absorber el par de reacción del volante.

	\item \textbf{Un plazo acotado.} El trabajo debía cerrarse antes de finales de agosto de 2026, lo
	que fijó el tamaño del conjunto de datos en cien demostraciones y la comparación en dos brazos.

\end{itemize}

Ninguna de ellas se oculta en el análisis: el capítulo \ref{Evaluación} dedica una sección a las
amenazas a la validez que derivan de ellas.


\section{Normativa asociada}

Todo el trabajo se desarrolla en simulación, de modo que ninguna norma le es directamente exigible.
Pero el sistema que estudia —un robot autónomo que manipula válvulas en una planta de proceso— está
regulado en todos sus aspectos, y un dictamen sobre su viabilidad que ignorase ese marco estaría
incompleto. Se recoge aquí la normativa que afectaría a un despliegue, organizada en cuatro bloques,
señalando en cada uno qué implicaría para un sistema basado en aprendizaje. La tabla
\ref{tab:normativa} la resume.

\textbf{Seguridad de máquinas y robótica.} Un robot es una máquina, y el marco legal europeo para las
máquinas es la Directiva 2006/42/CE, que el Reglamento (UE) 2023/1230 sustituye a partir de enero de
2027. El nuevo reglamento es relevante para este trabajo por una razón concreta: considera de forma
explícita los componentes de seguridad con comportamiento total o parcialmente autoevolutivo basado
en aprendizaje automático, y les exige evaluación de conformidad por un organismo notificado. Por
debajo del marco legal, ISO 12100 fija la metodología de evaluación del riesgo; ISO 10218-1 y -2
recogen los requisitos del robot industrial y de su integración; ISO/TS 15066 los límites de
contacto para robots que comparten espacio con personas, que un humanoide en una planta comparte
por definición; e ISO 13849-1 el diseño de las partes del sistema de mando relacionadas con la
seguridad.

\textbf{Atmósferas potencialmente explosivas.} Es el bloque determinante en el contexto del petróleo
y el gas. La Directiva 2014/34/UE, transpuesta por el Real Decreto 144/2016, regula los equipos
destinados a zonas clasificadas, y la Directiva 1999/92/CE, transpuesta por el Real Decreto
681/2003, la clasificación de esas zonas y la protección de los trabajadores; la serie IEC 60079
contiene los requisitos técnicos del material eléctrico. Un humanoide —baterías de litio, motores,
electrónica de potencia— es, desde este punto de vista, un equipo eléctrico que habría que certificar
para la zona en la que opere, y esa certificación es, hoy, el obstáculo más serio para llevar un
sistema como este a una zona clasificada. Un despliegue realista empezaría en áreas no clasificadas
o en zona 2.

\textbf{Seguridad funcional.} IEC 61508 establece los requisitos de los sistemas eléctricos,
electrónicos y programables relacionados con la seguridad, e IEC 61511 los particulariza para los
sistemas instrumentados de seguridad de la industria de proceso, que son los que protegen la planta
en la que el robot actuaría. El informe técnico ISO/IEC TR 5469 aborda la relación entre seguridad
funcional y sistemas de inteligencia artificial, y su conclusión práctica es la que cualquier
arquitectura de despliegue debería adoptar: la política aprendida no puede estar en el camino de
seguridad. Los límites de par, la parada de emergencia y la vigilancia de la posición del robot
deben garantizarlos capas deterministas, y la política solo debe decidir cómo se ejecuta la tarea
dentro de ellos.

\textbf{Inteligencia artificial.} El Reglamento (UE) 2024/1689 clasifica como de alto riesgo los
sistemas de inteligencia artificial que son componentes de seguridad de productos cubiertos por la
legislación de armonización, entre ellos las máquinas. Un sistema como el estudiado quedaría en esa
categoría, con las obligaciones que conlleva: gestión de riesgos, gobernanza de los datos de
entrenamiento, documentación técnica, registro de la actividad, supervisión humana y requisitos de
exactitud y robustez. La documentación del conjunto de datos y del protocolo de evaluación que esta
memoria recoge, y que en un trabajo académico es una cuestión de rigor, sería en un despliegue una
obligación legal.

\begin{table*}[ht]
	\centering
	\caption{Normativa aplicable a un despliegue del sistema}
	\label{tab:normativa}
	\makebox[\textwidth]{
		\begin{tabular}{|p{2.6cm}|p{3.7cm}|p{8cm}|}
			\cline{1-3}
			Bloque & Referencia & Ámbito \\ \hline
			Máquinas y robótica & Reglamento (UE) 2023/1230 & Máquinas; sustituye a la Directiva 2006/42/CE e incluye las funciones de seguridad basadas en aprendizaje automático \\
			 & ISO 12100 & Principios generales de diseño y evaluación del riesgo \\
			 & ISO 10218-1/-2 & Robots industriales: requisitos del robot y de su integración \\
			 & ISO/TS 15066 & Robots colaborativos: límites de contacto entre robot y persona \\
			 & ISO 13849-1 & Partes de los sistemas de mando relativas a la seguridad \\ \hline
			Atmósferas explosivas & Directiva 2014/34/UE & Equipos para atmósferas potencialmente explosivas (RD 144/2016) \\
			 & Directiva 1999/92/CE & Clasificación de zonas y protección de los trabajadores (RD 681/2003) \\
			 & IEC 60079 (serie) & Requisitos técnicos del material eléctrico para atmósferas explosivas \\ \hline
			Seguridad funcional & IEC 61508 & Sistemas eléctricos, electrónicos y programables relacionados con la seguridad \\
			 & IEC 61511 & Sistemas instrumentados de seguridad en la industria de proceso \\
			 & ISO/IEC TR 5469 & Seguridad funcional y sistemas de inteligencia artificial \\ \hline
			Inteligencia artificial & Reglamento (UE) 2024/1689 & Reglamento europeo de inteligencia artificial: sistemas de alto riesgo \\ \hline
		\end{tabular}
	}
\end{table*}

Por último, y en otro plano, las referencias bibliográficas de esta memoria siguen la norma ISO 690,
tal como establece la plantilla de la titulación.
